\documentclass{report}
\usepackage{amsmath,amssymb}
\setlength{\parindent}{0mm}
\setlength{\parskip}{1em}
\begin{document}
\begin{center}
\rule{6in}{1pt} \
{\large Alexander Toth \\
{\bf Anderson Acceleration for Black-Box Code Coupling in Nuclear Reactor Simulation}}

Department of Mathematics \\ Box 8205 \\ NC State University \\ Raleigh \\ NC 2695-8205
\\
{\tt artoth@ncsu.edu}\\
Tim Kelley\\
Roger Pawlowski\end{center}

Picard iteration is a standard method for solving coupled multi-physics
problems. In this, individual physical systems are repeatedly solved in
some sequence with updated coupling parameters transferred between
systems as they are obtained. Picard iteration is primarily attractive
due to its simplicity of implementation and minimal requirements on the
single-physics application codes. However, this method comes with several
drawbacks, namely relatively slow convergence and poor robustness.
Anderson acceleration is an alternative method for solving fixed-point
problems which maintains a depth of previous iterate information in order
to compute a new iterate as a linear combination of previous evaluations
of the fixed-point map. In this presentation, we examine the potential
for Anderson acceleration to improve upon the drawbacks of Picard
iteration in the context of coupled multi-physics problems in nuclear
reactor simulation. We specifically consider the Tiamat code coupling
being developed as part of the Consortium for Advanced Simulation of
LWRs. Tiamat couples the Bison fuel performance code with the MPACT
neutronics and COBRA-TF thermal hydraulics codes to provide a tool for
pellet-cladding interaction analysis. We will describe how Anderson
acceleration has been integrated into this code coupling by posing the
coupled system as a fixed-point problem in terms of coupling parameters,
and then examine the performance gains obtained from utilizing Anderson
acceleration for this coupling.


\end{document}
